% =============================================================
% ESTIMATING LATVIAN GDP, 1937-1995
% =============================================================
\section{Estimating Latvian GDP, 1937--1995}
\label{sec:method}

\subsection{The output approach and the reliability of Soviet statistics}

For the Soviet period only the production side of national income can be reconstructed. The monetary aggregates that an expenditure or income estimate would require are the figures that cannot be trusted: the Material Product System measured net material product at administered prices, excluded most services, double-counted intermediate transactions in its gross-output measures and, above all, understated inflation, so that official constant-price growth was too high; in Maddison's verdict, the official series `overstated growth, and understated inflation' \parencite[307]{Maddison1998b}. The Western response, begun by Bergson in the 1950s and institutionalised by the CIA from the 1960s, was to rebuild Soviet GNP from the bottom up, re-pricing output at factor cost, imputing the missing services and replacing the official value indices with indices built from physical output series weighted by value-added \parencite{bergson1950,bergson1961,cia1982a}. The CIA estimates were attacked from both sides, by critics who held that they overstated Soviet growth by ignoring deteriorating quality, over-reporting and hidden inflation, and faulted even the physical series as `open' statistics compiled for external consumption, and by a minority who contended that military output, housing and services were under-measured \parencite{assessingCIA1988,becker1994,schroeder1995}; Khanin's alternative statistics, as \textcite[144, 159]{harrison1993} showed, coincided with the CIA's for 1950--87. The assessments of Becker and Schroeder and the review by Maddison conclude that the CIA series contain errors in both directions, that most of the criticism `was not well founded', and that the physical-output method remains the best basis for measuring Soviet growth; the CIA's estimates, in their final version at 1982 prices \parencite{cia1990}, are the basis of Maddison's Soviet series, and Vonyó and Markevich describe the Western revisions as still widely acknowledged as the best available \parencites[309]{Maddison1998b}[280]{vonyoinmorys2021}. Ponomarenko's archival reconstruction of Russian national accounts reaches the same judgement from inside the Soviet statistical system: the physical output data `are apparently of high quality and can be used', whereas the official value indices were inflated by the pricing of new products \parencite[16--18, 73--4]{ponomarenko2002}. The literature on the Soviet economy has grown rapidly since the opening of the archives \parencite{zhuravskaya2024}, but its quantitative core for the post-war decades still rests on physical output series of this kind, and the non-Russian republics remain almost entirely without modern national accounts, studies of Soviet growth being conducted at the level of the Union or of Russia \parencites[646, 670--1]{kukic_soviet}{voskoboynikov2021}.

The series presented here follows this tradition and the procedure used for Lithuania and Estonia \parencite{klimantas2026soviet,klimantas2025estonia}. A benchmark estimate of GDP is extrapolated with constant-price indices of sectoral value-added, each sector weighted by its share of GDP in the benchmark year:
\[
    \mathit{GDP}_{t} = \mathit{GDP}_{\mathrm{bench}} \times \sum_{i=1}^{n} w_{i}\, I_{it},
\]
where $w_i$ is sector $i$'s share of benchmark GDP and $I_{it}$ its volume index in year $t$. The indices track gross output, inputs or employment and are combined with value-added weights, what Thomas and Feinstein call the indicator approach, as distinct from single deflation, although like it the procedure holds the ratio of value-added to output constant; it is the procedure that Fabricant, as they quote him, called the `next best' where double deflation is impossible because intermediate inputs cannot be observed \parencite[427--8, 444]{thomasfeinstein2003}, and the CIA itself double-deflated only agriculture \parencite[45]{cia1982a}. It matters, therefore, which Soviet data are used and which avoided. The distortions documented in the literature were concentrated in value aggregates, in the national income and gross output indices at `comparable prices', in the reporting of hidden price increases as real growth and in the double counting of the material product system \parencites[307--8]{Maddison1998b}[6]{vonyo2017}[165]{mertelsmann_2003_there}; Khanin's own classification, as Harrison summarises it, is that `the ease of distortion remained greater for value series than for physical volumes', and that distortion was concentrated in heterogeneous products such as machinery and construction, while series for grain, milk, cement, electricity or freight are comparatively reliable \parencite[144, 152]{harrison1993}. The two diverge widely: the official Soviet index of Latvian industrial output at constant prices rose nineteen-fold between 1940 and 1980 \parencite[361]{years_of_dependence}, more than twice the increase recorded by the physical-product index compiled here (Table~\ref{tab:sectoral}). The official index also values its 1940 base at an arbitrary lat--rouble parity that undervalued the last pre-Soviet year several-fold (Misiunas and Taagepera suspected in addition that only the five months of Soviet rule were counted; Mertelsmann finds for Estonia that the whole year was); its claim that industry had surpassed 1940 by 1947 is contradicted by the physical series \parencites[109]{years_of_dependence}[167]{mertelsmann_2003_there}[234]{krumins_2017_sovetizacijas}. The indices below rest almost entirely on physical quantities, on employment and on inputs; the one place where value data at comparable prices enter, trade turnover, which also scales the finance index, is flagged and its bias bounded in the sensitivity analysis.

\subsection{The 1990 benchmark}

The Soviet-period series is anchored to Bolotin's estimate of the GDP of the Latvian SSR in 1990, which expresses the accounts of all fifteen republics in international dollars \parencite[181--92]{bolotin1992}. The year 1990 is one of the first for which republican accounts on an SNA-type definition were compiled \parencite[43--4]{gulans_2015_rupnieciba} and the first with purchasing-power parities; Bolotin converted the republican GDP of that year into dollars with parities estimated, as he put it, `in practically full conformity' with the International Comparison Project. The parities, by sector of origin, were 0.40 dollars per rouble in agriculture, 1.80 in industry, 1.10 in construction, 2.80 in transport, 1.45 in trade and 3.30 in other services, or 1.82 for GDP as a whole \parencite[181, 184]{bolotin1992}. For Latvia the result is \lv{1990}~\intlnineteen{} per head, between Lithuania's 7,475 and Estonia's 8,425. Table~\ref{tab:benchmarks} shows the sectoral composition in roubles, with services other than trade and transport allocated by the 1987 Soviet input--output table \parencite{wbinputoutput1987}, beside the 1935 benchmark of \textcite{norkuslatvia1935} on which the interwar series rests. The rouble shares are domestic-price weights, like the lats shares of 1935, and are the weights with which the volume indices are combined; the dollar conversion applies to the total alone. They are at the market prices of 1990, so that product taxes and subsidies are carried within the sectors that bore them; no re-pricing to adjusted factor cost was attempted, since no republican data on the incidence of net product taxes by branch exist, and the only socialist economies for which the World Bank reports the ratio of product taxes less subsidies to value-added, Bulgaria and Ukraine in 1980--90, show about 0.5 per cent \parencite{wdi2024}.\footnote{The turnover tax, the principal product tax, was levied on final consumption goods and accrued to the republic where the finished good was made, as with cotton fabric woven in the Baltics from cheaply procured Uzbek cotton \parencite[20]{yasin_1993_economic}.}

\authorcheck{\textbf{Author check (Table~\ref{tab:benchmarks}).} The 1990 industry value, 3,190.8 million roubles, does not follow from Bolotin: his 5.2 billion dollars at the industry parity of 1.80 give 2,888.9 million roubles, while every other sector reproduces exactly \parencite[184]{bolotin1992}. The dissertation's benchmark table carries the same uplift (+301.9 for Latvia, +307.2 for Lithuania, +259.0 for Estonia). State the adjustment here and in the source note of Table~\ref{tab:benchmarks}, or revert to Bolotin's figure.}

% tab:benchmarks
\begin{table}[htbp]
\centering\small
\caption{Benchmark GDP of Latvia by sector of origin, 1990 and 1935, current prices}
\label{tab:benchmarks}
\begin{threeparttable}
\begin{tabular}{lrrrr}
\toprule
& \multicolumn{2}{c}{1990} & \multicolumn{2}{c}{1935} \\
\cmidrule(lr){2-3}\cmidrule(lr){4-5}
Sector & m.\,roubles & \% & m.\,lats & \% \\
\midrule
Agriculture, forestry and fishing & 2,000.0 & 17.6 & 383.4 & 36.1 \\
Industry & 3,190.8 & 28.0 & 216.5 & 20.4 \\
Construction & 818.2 & 7.2 & 23.6 & 2.2 \\
Transport and communications & 964.3 & 8.5 & 53.2 & 5.0 \\
Trade, accommodation and catering & 1,448.3 & 12.7 &  &  \\
Finance and insurance & 234.1 & 2.1 &  &  \\
Housing & 689.2 & 6.1 &  &  \\
\quad \textit{Commercial services, 1935 total} & & & \textit{233.6} & \textit{22.0} \\
Education and culture & 580.2 & 5.1 &  &  \\
Health & 357.1 & 3.1 &  &  \\
Science & 203.4 & 1.8 &  &  \\
Public administration & 114.8 & 1.0 &  &  \\
Defence & 626.9 & 5.5 &  &  \\
\quad \textit{Public services, 1935 total} & & & \textit{103.5} & \textit{9.7} \\
Other services & 163.9 & 1.4 & 48.0 & 4.5 \\
\midrule
Gross value-added & 11,391.2 & 100.0 & 1,061.8 & 100.0 \\
Taxes less subsidies on products & --- & & 77.8 & \\
GDP at market prices & 11,391.2 & & 1,139.6 & \\
GDP per capita, national currency & 4,277.3 & & 583.4 & \\
\textit{GDP per capita, \intlnineteen{}} & \textit{7,965} & & \textit{2,776} & \\
\bottomrule
\end{tabular}
\begin{tablenotes}[para,flushleft]
\footnotesize
\textit{Notes:} the 1990 benchmark is at market prices: taxes less subsidies on products are carried within the sectors that bore them and are not reported separately (shown as `---'). The 1935 benchmark reports commercial and public services as two aggregates, shown in italics. Population: mid-year. \textit{Sources:} 1990: \textcite[184--6]{bolotin1992}, converted back to roubles with Bolotin's sectoral purchasing-power parities, services other than trade and transport split with the 1987 Soviet input--output table \parencite{wbinputoutput1987}; 1935: \textcite[263--6, 287, 312]{norkuslatvia1935}; per capita figures in \intlnineteen{}: see online appendix.
\end{tablenotes}
\end{threeparttable}
\end{table}


Three features of the benchmark need stating, because the reconstruction stands or falls with it. First, its level is conservative. Bolotin's figure for the Soviet Union as a whole, 6,275~\intlnineteen{} per head, lies below Maddison's own estimate for 1987, and the bilateral Austria--USSR comparison on which it draws placed the Union at 32 per cent of the United States where the CIA had placed it at 42 per cent \parencite[318--19]{Maddison1998b}. Maddison nevertheless built the Maddison Project's 1990 levels for the republics from Bolotin's republic-to-Union ratios of 1991 rather than 1990, ratios inflated by the first year of transition, and then carried them backwards with net material product growth \parencite[181--2]{maddisonworldeconomy}; the consequences for Latvia are shown in Table~\ref{tab:previous} below. Second, the parities are Union-wide, so that the conversion assumes that relative prices in Latvia were those of the Soviet Union as a whole. For the material sectors this is a reasonable approximation in a planned economy with administered prices; for services, where the parity of 3.30 dollars per rouble reflects the very low rouble wages of the non-market sector, it is less certain, and the resulting share of services in Latvian GDP at rouble prices, 47 per cent, is above the 34 per cent implied by the World Bank's 1990 sectoral data \parencite{wbstatisticalhandbook1994}; the sensitivity of the series to this allocation is examined below. Third, there is no alternative.\footnote{Earlier claims ranged widely: PlanEcon's CIA-based estimate for Estonia in 1989 was 6,240 dollars against 5,000 for the Union \parencite[166]{kukk2005}.} Ponomarenko's accounts cover Russia alone and make no dollar conversion; the World Bank's GDP levels for 1980--91 \parencites{wbtransitionlatv}[32--3]{worldtables1995} were not continued; the UNECE's `real GDP' indices for 1980--90 were taken from net material product \parencite[249]{UNECE1995}; the CIA's 1982 benchmark cannot be decomposed by republic, and the 1987 input--output table understates services too severely to serve as a benchmark \parencite{wbinputoutput1987}.

The choice of 1990 means that the series is extrapolated backwards for more than fifty years from end-of-period weights. The practical reason is that 1990 is the only year for which republican accounts on an SNA definition and purchasing-power parities exist together. Two further considerations favour it: the reporting of 1989--90, at the height of glasnost, was more transparent than that of any earlier Soviet year, the manipulation of earlier returns having been publicly acknowledged in 1988 by the deputy head of the Lithuanian statistical committee \parencite{adlys1988}; and 1990 is a year in which the branches built up after the war, chemicals and pharmaceuticals at Olaine, electrical engineering or microbus production, can be observed and weighted at all \parencites[278]{krumins_2017_tautsaimnieciba_1950}[303, 305]{krumins_2017_tautsaimnieciba_1960}. The cost is that half a century is valued at end-period relative prices, which, as the sensitivity analysis explains, tends to understate growth; the short answer is that the consequence can be tested directly, because the Soviet-period series meets an independent series with 1935 weights in 1937--40.

\subsection{Volume indices and their sources}

Table~\ref{tab:sources} summarises the indicators and sources for each sector; the online appendix documents every series year by year. The Soviet-period indices are built from the statistical yearbooks of the Latvian SSR, which from 1957 give annual physical output from the mid-1950s and, unusually among the Baltic republics, cover the early post-war years, though incompletely; from the statistical bulletins published by the German civil administration for 1942--4; from the USSR compendia on domestic trade and on the national economy; from the World Bank's studies of 1993--6; and from archival material, notably the returns of the Latvian statistical administration on sown areas for 1945--55 and the transport returns of 1945--50 held in the Russian State Archive of the Economy. Three general points precede the detail. First, archival material is used to fill the early post-war years that the yearbooks cover incompletely, not merely to decorate the series; but, unlike for Lithuania and Estonia, no archival compilation of late-1980s output values at the product level could be found for Latvia, so that the 129 product series are consolidated into 22 minor branches by averaging and the minor branches weighted with the 1987 input--output study and the World Bank's 1993 valuation \parencites[82--5]{wbinputoutput1987}[289]{wbtransitionlatv}. Second, the series cover the civilian economy as reported to the republican statistical office. Military production, organised through all-Union enterprises and never reported separately, is not captured, and Latvia hosted more of it than its neighbours: the VEF electrotechnical works, the Riga radio and instrument plants and the shipyards all carried `special' production for the Soviet armed forces \parencites[59--61]{upmalis_2015_latvijas}[74--5]{blaubergs_2015_vef}[103--4]{wbtransitionlatv}; the size of the Soviet garrison is likewise unknown between 1941 and 1989 and is treated as a source of bias below. Third, the year 1944 proved beyond reconstruction from Latvian sources, so that sectoral movements for 1943--5 are taken from the ratios observed in Lithuania, where documentation is better.

% tab:sources
\begin{table}[htbp]\centering\scriptsize\setlength{\tabcolsep}{3pt}
\caption{Indicators and sources of the Latvian Soviet-period series, 1937--95}
\label{tab:sources}
\begin{threeparttable}
\begin{tabular}{p{2.9cm}p{11.3cm}c}
\toprule
Sector & Indicator and main sources & Appendix \\
\midrule
Crops & Physical output of 9--10 crop categories, 1937--95; 1990 value-added weights from the World Bank's 1993 study, grains split with the 1995 Eurostat agricultural accounts; sown areas for 1945--55 from unpublished returns & A.1 \\
Livestock & Physical output of 7--16 animal products; 1941--9 yields obtained from animal numbers and the yields per animal of 1940, 1945 and 1950 & A.1 \\
Forestry and fishing & Timber cut (cubic metres) and fish catch (tonnes); 1941, 1971 and 1974--6 catches interpolated & A.1 \\
Industry & 129 physical product series in 11 branches (Table~\ref{tab:products}); product weights from the 1987 input--output study and the 1993 World Bank study, products consolidated into 22 minor branches by index averages because no archival late-1980s output valuation survives for Latvia; branch weights from the 1990 benchmark & A.2 \\
Construction & Weighted index of material inputs into construction, 1987 input--output weights, following the CIA procedure & A.3 \\
Transport and communications & Freight tonne-kilometres and passenger-kilometres, all modes; postal, telegraph and telephone traffic and radio and television sets; 1945--50 from the Russian State Archive of the Economy; 1941--4 assumed to move with the rest of the economy; 1990s from sectoral employment & A.4 \\
Trade, accommodation and catering & Soviet retail and wholesale turnover indices at comparable prices (the only value series used), chained across base-year changes; early post-war years from the USSR domestic-trade compendia & A.5 \\
Finance, insurance and housing & Finance and insurance: employment by branch, adjusted for productivity with the ratio of trade turnover to trade employment, as the CIA did; housing: housing stock derived from construction and population & A.5 \\
Public administration, health, education, science & Employment by branch from Soviet yearbooks, the Ostland statistical bulletins and World Bank studies; 1941--4, 1946--9 and 1951--4 gaps interpolated & A.6 \\
Defence & Soviet garrison of 30,000 in 1940 and 80,000 in 1990 plus the republican forces and police of 1940; 1941--89 interpolated between the 1940 and 1990 totals in step with public-administration employment; 1991--5 from the withdrawal schedule and the new Latvian forces & A.6 \\
Other services & Employment & A.6 \\
Population & Year-end estimates and censuses of 1935, 1959, 1970, 1979 and 1989; wartime and deportation losses from the demographic literature; Soviet garrison added in 1940--1; present-day borders from 1945 & A.7 \\
\bottomrule
\end{tabular}
\begin{tablenotes}[para,flushleft]
\footnotesize
\textit{Notes:} the appendix column refers to the sections of the online appendix in which every series is documented year by year with its sources: the Soviet Latvian statistical yearbooks (\textit{Latvijas PSR tautas saimniec\={\i}ba}, 1957--90), the \textit{Statistische Berichte f\"{u}r das Ostland} (1942--4) and \textit{Ostland in Zahlen} (1942), the USSR domestic-trade compendia (\textit{Sovetskaya torgovlya}, 1964 and 1989), the World Bank's \textit{Latvia: the transition to a market economy} (1993) and \textit{Statistical handbook: states of the former USSR} (1994--6), the Russian State Archive of the Economy (RGAE) and the Latvian State Archives. \textit{Sources:} see online appendix.
\end{tablenotes}
\end{threeparttable}
\end{table}


% tab:products
\begin{table}[htbp]
\centering\footnotesize
\caption{Industrial product series compiled for the Soviet-period indices, 1937--95}
\label{tab:products}
\begin{threeparttable}
\begin{tabular}{lrrr}
\toprule
Branch & Latvia & Lithuania & Estonia \\
\midrule
Electricity & 1 & 1 & 1 \\
Petroleum & 0 & 1 & 3 \\
Other fuels & 3 & 1 & 2 \\
Metals & 4 & 7 & 4 \\
Chemicals & 8 & 10 & 11 \\
Machine-building and metal-working & 30 & 36 & 21 \\
Wood and paper & 9 & 14 & 10 \\
Construction materials & 22 & 15 & 8 \\
Light industry & 18 & 18 & 9 \\
Food industry & 25 & 36 & 25 \\
Other & 9 & 4 & 7 \\
\midrule
Total & 129 & 143 & 101 \\
\bottomrule
\end{tabular}
\begin{tablenotes}[para,flushleft]
\footnotesize
\textit{Sources:} this paper; \textcite{klimantas2026soviet}; \textcite{klimantas2025estonia}.
\end{tablenotes}
\end{threeparttable}
\end{table}


\textit{Agriculture, forestry and fishing.} The index combines physical output of nine to ten crop categories and seven to sixteen animal products with timber removals and the fish catch, each sub-index weighted by its 1990 value-added, the crop and livestock weights coming from the World Bank's 1993 study and the split within grains from the Eurostat agricultural accounts \parencite[286]{wbtransitionlatv}. Early post-war output comes from unpublished returns on sown areas; yields for 1941--9 are obtained from animal numbers and the yields per animal of 1940, 1945 and 1950; disaggregated grain data for 1947--52 are not available, so that the grain categories move with the total.

\textit{Industry.} Physical output series could be assembled for 129 products in ten of the eleven branches of the benchmark, electricity, fuels, metals, chemicals, machine-building and metal-working, wood and paper, building materials, light industry, food and a residual; the eleventh, petroleum refining, did not exist in Latvia (Table~\ref{tab:products}). The fuel index for 1980--90 is taken from the World Bank's study and extended with peat and natural gas; where a product series begins or ends with zero, production is assumed absent in the adjacent years. The branches are combined with the branch structure of the 1990 benchmark. For the interwar years the published manufacturing volume index for 1930--9 is used and the earlier years rest on employment adjusted for real wages, as in the Lithuanian and Estonian series \parencite[39, 45]{ostlandinzahlen}.

\textit{Construction.} The Soviet-period index is a weighted index of the material inputs into construction, with weights from the 1987 input--output table, which is the CIA's procedure \parencite[66--89]{wbinputoutput1987}.

\textit{Transport and communications.} Freight tonne-kilometres and passenger-kilometres across all modes are combined with a communications index built from postal and telegraph traffic, long-distance calls and radio and television sets in use, with weights from the 1987 input--output study. The years 1945--50 come from the Russian State Archive of the Economy; for 1941--4 no transport data exist and the sector is assumed to move with the rest of the economy; after 1990 the index follows sectoral employment.

\textit{Trade, finance and housing.} Under planning industrial products circulated by allocation rather than through markets, so that the marketed-share approach of the interwar series cannot be used. The Soviet-period trade index instead chains the retail and wholesale turnover indices at comparable prices compiled by the Soviet statistical offices, the early post-war years coming from the USSR domestic-trade compendia; these are the only value series in the reconstruction. Finance and insurance follow employment, as in the CIA's estimates, which used employment as a proxy for output in such sectors \parencites[319]{Maddison1998b}[49]{cia1982a}; unlike the CIA, which found no evidence on which to base a productivity allowance, the present series adjusts employment for productivity with the ratio of trade turnover to trade employment. Housing follows the housing stock, as in the CIA's index \parencite[46]{cia1982a}, here derived from construction and population.

\textit{Public and other services.} Public administration, health, education and science follow employment by branch, as in the CIA's accounts, with gaps in 1941--4, 1946--9 and 1951--4 interpolated. Defence is the least secure component. The Soviet garrison admitted under the mutual-assistance pact of October 1939 numbered about 30,000 \parencite[15--16]{years_of_dependence}, and some 80,000 Soviet troops were stationed in Latvia in 1990, far more in proportion to population than in Lithuania or Estonia; between those dates the garrison is interpolated between the 1940 and 1990 totals in step with public-administration employment, and after 1990 the index follows the withdrawal schedule and the formation of the Latvian forces.

\textit{The war years.} Latvia was under German civil administration from 1941 to 1944, and the \textit{Reichskommissariat Ostland} published statistics, the \textit{Statistische Berichte für das Ostland} of 1942--4 and \textit{Ostland in Zahlen}, that cover crops, livestock, industrial output and employment but little else. Where these fail, the service sectors are interpolated between 1940 and the first post-war observation, and the movement between 1943 and 1945 is taken from the Lithuanian ratios. The estimates for 1941--5 are therefore `rough' in Feinstein's sense, with a margin of 15--25 per cent, and the 1944 trough is an order of magnitude rather than a measurement \parencite[21]{feinstein1972}. Because the Soviet-period series is extrapolated backwards through the war, an error in the wartime linkage would shift its level in 1937--40, so that the overlap tests the linkage together with the rest of the extrapolation; the spliced series, which follows the interwar series to 1940, is affected only through the 1940 average.

\subsection{Population and borders}

Latvia's borders changed once, when the eastern Abrene district was transferred to the Russian republic in 1944. The series to 1944 refer to the borders of the republic and the series from 1945 to present-day borders. This is the convention of the Maddison Project, which treats income per head as unaffected by a territorial change that moves less than a third of the population \parencite[629]{maddison_project_paper2014}; the transferred district held a few per cent of the pre-war population at most \parencites[74]{years_of_dependence}[129--30]{hiden_salmon}, and was rural and poorer than the national average, so that the break in income per head is small and no adjustment has been made. The population series rests on the censuses of 1920, 1925, 1930, 1935, 1959, 1970, 1979 and 1989 and the year-end estimates of the Latvian statistical office and the `Baltic 100' database, with mid-year figures as the average of consecutive year-ends; wartime losses, the loss of some 120,000 ethnic Latvians as refugees and repatriates in 1939--49 \parencite[86]{king_1965_economic}, the Holocaust and the deportations are taken from the demographic literature, and the Soviet garrison of 1940--1 is added to the resident population. Between 1939 and 1945 the resident population fell by \lvPOPchangeWARtwo{} per cent, of which a small part, of the order of two points, reflects the border change (Table~\ref{tab:series}).

\subsection{The interwar series and its extension to 1919 and 1940}

The interwar series is that of \textcite{Klimantas2024}, which extrapolates the 1935 benchmark of \textcite{norkuslatvia1935} to 1920--39 with sectoral volume indices built from the publications of the Latvian statistical office;\footnote{Bukowski, Kowalski and Wroński object that this extrapolation holds relative prices fixed at their 1935 structure although agricultural--industrial relative prices fluctuated between the wars \parencite[6]{bukowski_2026_economic}; the objection is the index-number problem discussed under sensitivity below.} the benchmark itself was converted into international dollars by price comparisons with Sweden and Lithuania, of the kind used for Lithuania in 1937 and Estonia in 1938. Two years are added here. For 1919, the first year of the republic's existence, the crop, livestock, railway and employment series were carried back with the data for 1919 where they exist and with 1920 ratios where they do not; forestry, fishing, finance and the household services follow the population or the sectors to which they are tied. For 1940, the year of the Soviet occupation, the 1940 yearbook, the Ostland publications and the first Soviet returns provide output and employment for most sectors, imports are assumed equal to their deflated 1939 value, and defence includes the Soviet garrison. The resulting figure for 1919, \lv{1919}~\intlnineteen{}, is 37 per cent below the 1913 benchmark of Norkus and Markevičiūtė, and places Latvia, with Lithuania, among the most devastated economies in Europe at the end of the First World War; contemporaries and later historians judged Latvia perhaps the most devastated of all \parencites[323]{norkusdvides}[44]{krumins_2024_makroekonomiska}.

\subsection{Aggregation, splicing and the overlaps}

Each of the two new series is compiled in constant benchmark-year prices and expressed in \intlnineteen{} through its benchmark. The interwar series covers 1919--40 and the Soviet-period series 1937--95, so that the two overlap for four years; the Soviet-period series in turn overlaps for six years with a modern series obtained by backcasting the Maddison Project's 2011 benchmark with the World Bank's constant-price growth rates of GDP per capita, which are used in preference to the Maddison Project's own 1990--2010 growth for reasons given below \parencite{mpd2023,wdi2024}. Table~\ref{tab:overlaps} reports the overlaps.

% tab:overlaps
\begin{table}[htbp]
\centering\small
\caption{Overlapping estimates of Latvian GDP per capita at the two splicing points}
\label{tab:overlaps}
\begin{threeparttable}
\begin{tabular}{lrrcc}
\toprule
Year & Earlier-benchmark series & Later-benchmark series & Difference (\%) & Feinstein grade \\
& (\intlnineteen{}) & (\intlnineteen{}) & & \\
\midrule
\multicolumn{5}{l}{\textit{(a) Interwar series (1935 benchmark) and Soviet-period series (1990 benchmark)}} \\
1937 & 3,031 & 3,082 & +1.7 & Firm \\
1938 & 3,099 & 3,198 & +3.2 & Firm \\
1939 & 3,108 & 3,049 & $-$1.9 & Firm \\
1940$^{\dagger}$ & 3,006 & 2,985 & $-$0.7 & Firm \\
\midrule
\multicolumn{5}{l}{\textit{(b) Soviet-period series (1990 benchmark) and adjusted MPD series (2011 benchmark)}} \\
1990 & 7,965 & 8,196 & $-$2.9 & Firm \\
1991$^{\dagger}$ & 7,385 & 7,205 & +2.4 & Firm \\
1992 & 6,280 & 4,959 & +21.0 & Rough \\
1993 & 5,625 & 4,806 & +14.6 & Good \\
1994 & 5,636 & 4,994 & +11.4 & Good \\
1995 & 5,710 & 5,018 & +12.1 & Good \\
\bottomrule
\end{tabular}
\begin{tablenotes}[para,flushleft]
\footnotesize
\textit{Notes:} Difference = (later $-$ earlier)/earlier in panel (a) and (earlier $-$ later)/earlier in panel (b), so that a positive number always means that the 1990-based series is the higher one. Feinstein's grades: firm, under 5 per cent; good, 5 to under 15; rough, 15 to under 25; conjectural, 25 and over \parencite[21]{feinstein1972}. $^{\dagger}$~Splice year: the final series takes the arithmetic average of the two estimates. \textit{Sources:} see online appendix; adjusted MPD series from \textcite{mpd2023} with 1990--2010 growth replaced by World Bank constant-price growth backcast from the 2011 benchmark \parencite{wdi2024}.
\end{tablenotes}
\end{threeparttable}
\end{table}


In 1937--40 the two series differ by $+$1.7, $+$3.2, $-$1.9 and $-$0.7 per cent, all within the 5 per cent margin that Feinstein reserved for `firm' estimates, and not of one sign. The test is demanding. The two series share no sources, the one resting on the republic's own statistics and the other on Soviet yearbooks and post-war archives; no weights, one using the structure of 1935 and the other that of 1990; and no benchmark, one a price comparison with Sweden and Lithuania, the other a purchasing-power-parity study of the Soviet republics fifty-five years later. That the Soviet-period series, extrapolated backwards over half a century of war, collectivisation and forced industrialisation, arrives within 3.5 per cent of an independently measured level in every year, and within 2 per cent in three of the four, is the strongest evidence available that the net error of the extrapolation is small. It cannot exclude errors that happen to offset, but the two index-number biases discussed below, late base-year weights and fixed value-added weights, would both understate growth and so would both place the backcast 1937--40 level above the interwar series; a net understatement of as little as 0.1 percentage point a year over 1940--90 would place it 5 per cent above in every overlap year, which is not observed. \textcite[323--6]{broad2007} regarded time-series projections as consistent with direct benchmarks when they agree within 10 per cent; the Latvian overlap meets that standard in each of four consecutive years, as the Lithuanian and Estonian series do in 1939--40 and 1937--40 \parencite{klimantas2026soviet,klimantas2025estonia}. The two series are spliced at 1940 by taking the average of the two estimates for that year, which moves the 1940 figure by less than half a per cent.

The second overlap validates the benchmark. In 1990 the Soviet-period series stands 2.9 per cent below the modern series and in 1991 \ovd{1991} per cent above it: Bolotin's purchasing-power parities and the Maddison Project's 2011 ICP round, carried backwards with the national accounts, are two routes that share no data and arrive at the same level within 3 per cent. The series are spliced at 1991, the year of restored independence, by averaging. Thereafter the gap widens to \ovd{1992} per cent in 1992 before narrowing to \ovd{1995} per cent by 1995, because the modern series, built on the official national accounts, records the deeper fall in 1992, in line with the view that official data overstated the early-transition fall in output and missed the new private activities that were emerging unrecorded \parencites[7, 11]{imf_transition}[260]{kyn_1997_eastern}, while the Soviet-period series, valuing output at 1990 weights, records a shallower one; Hamilton and de Vries's recent reconstruction of the post-Soviet republics likewise extrapolates backwards from the most recent official levels \parencite[17]{hamilton_transition}. The two probably bracket the true fall; the modern series is preferred from 1992. It is anchored on the Maddison Project's 2011 figure for Latvia and carried back with World Bank growth because the Maddison Project's own path for 1990--2010 implies far slower growth than the national accounts record, so that its 1990 level, 15,806~\intltwenty{} or 9,916~\intlnineteen{}, stands \mpdgap{1990} per cent above Bolotin's independent estimate. That level is Maddison's, built on Bolotin's republic-to-Union ratios for 1991 rather than 1990 (Section~\ref{sec:gap}) and retained unchanged in the 2023 release; since the Maddison Project's path and the World Bank's share the 2011 anchor, the inflated 1990 level appears as slower growth over 1990--2010 in the Maddison Project's path, and its growth rates converge with the World Bank's only from 2011 \parencites[57]{klimantas2025thesis}{mpd2023}{wdi2024}.

\subsection{The final series and earlier estimates}

The 1913 figure is the benchmark of \textcite[603]{norkusmark2021}, taken as published; there are no estimates for 1914--18. Following \textcite[30]{bolt_zanden_newgdp2021}, the series is compiled in \intlnineteen{} and converted to \intltwenty{} at the ratio of 1.594 used by \textcite{mpd2023}, and both versions are reported, with the population, for 1913 and every year from 1919 to 2024 in Table~\ref{tab:series} of the appendix; Table~\ref{tab:previous} compares it with the earlier estimates at the years for which they exist.

% tab:previous
\begin{table}[htbp]
\centering\small
\caption{Latvian GDP per capita: the new series and earlier estimates, selected years}
\label{tab:previous}
\begin{threeparttable}
\begin{tabular}{lrrrr}
\toprule
Year & This paper & \multicolumn{3}{c}{Earlier estimates (difference from this paper, \%)} \\
\cmidrule(lr){3-5}
& (\intlnineteen{}) & Ros\'{e}s and Wolf (2010) & Norkus and Markevi\v{c}i\={u}t\.{e} (2021) & MPD 2023 \\
\midrule
1913 & 2,361 & --- & 2,361 (+0.0) & --- \\
1922 & 1,654 & 1,929 (+16.6) & 1,847 (+11.7) & --- \\
1929 & 2,351 & 2,798 (+19.0) & 2,347 ($-$0.2) & --- \\
1938 & 3,099 & 4,048 (+30.6) & 2,836 ($-$8.5) & --- \\
1973 & 8,523 & --- & --- & 7,846 ($-$7.9) \\
1980 & 8,867 & --- & --- & 8,619 ($-$2.8) \\
1985 & 9,284 & --- & --- & 8,954 ($-$3.6) \\
1989 & 8,772 & --- & --- & 9,825 (+12.0) \\
1990 & 7,965 & --- & --- & 9,916 (+24.5) \\
\bottomrule
\end{tabular}
\begin{tablenotes}[para,flushleft]
\footnotesize
\textit{Notes:} MPD figures are published in \intltwenty{} and converted at 1.594. The 1913 figure of this paper is that of Norkus and Markevi\v{c}i\={u}t\.{e}, taken as published. \textit{Sources:} \textcite[190]{roseswolf2010}; \textcite[603]{norkusmark2021}; \textcite{mpd2023}; this paper (online appendix).
\end{tablenotes}
\end{threeparttable}
\end{table}


For the interwar years the new series lies between the two indirect approximations, above Norkus and Markevičiūtė's and well below Rosés and Wolf's, whose 1938 figure would have made Latvia richer than Finland. For the Soviet period the comparison with the Maddison Project is instructive in a way that differs from the Lithuanian and Estonian cases. The Maddison Project's Latvian levels for 1973--85, carried back from a 1990 level with net material product growth, lie 3--8 per cent \textit{below} the new series, whereas its 1989 level lies \mpdgap{1989} per cent above it and its 1990 level \mpdgap{1990} per cent above: the Maddison Project records Latvian income per head rising by 5 per cent in 1989 and further in 1990, when every physical indicator compiled here was already falling; the EBRD's index of real GDP likewise rises by 6.8 per cent in 1989 and 2.9 per cent in 1990, whereas Krūmiņš's physical series put 1990 output below that of 1985 \parencites[269]{kyn_1997_eastern}[348--50]{krumins_2017_gorbacova}. The discrepancy is not in the benchmark alone but in the growth attributed to the late 1980s, which the net material product series overstated, by about 2 percentage points a year on United Nations estimates from around 1980 \parencite[289]{years_of_dependence}.

\subsection{Limitations and sensitivity}

Table~\ref{tab:sensitivity} lists the sources of bias, the sectors they affect and the maximum impact of each on the level of GDP, with the bounds derived as in the companion studies.

% tab:sensitivity
\begin{table}[htbp]
\centering\footnotesize
\caption{Sensitivity of the Latvian series to the main sources of bias}
\label{tab:sensitivity}
\begin{threeparttable}
\begin{tabular}{p{4.6cm}lp{3.8cm}rr}
\toprule
Source of bias & Years & Sectors affected & Bias & Max.\ impact \\
& & & (\%) & on GDP (\%) \\
\midrule
Incorrect service-sector evaluation & 1919--40 & Services & $\pm$10 & $\pm$3 \\
Fixed value-added weights (single deflation) & 1919--40 & All & $\pm$2 & $\pm$2 \\
Labour productivity and real-wage estimates & 1919--29 & Industry & $\pm$10 & $\pm$2 \\
Missing annual product-tax data & 1919--40 & Product taxes & $\pm$9 & $\pm$1 \\
No productivity adjustment in employment series & 1919--40 & Public services & $\pm$10 & $\pm$2 \\
Population movement as proxy & 1919--40 & Various service subsectors & $\pm$30 & $\pm$1 \\
\midrule
Incorrect service-sector evaluation & 1940--90 & Services (excl. transport) & $\pm$10 & $\pm$4 \\
Missing product-tax data & 1940--90 & Total & $-$0.5 & $-$0.5 \\
Fixed value-added weights (single deflation) & 1940--90 & All & $-$14 & $-$14 \\
Hidden inflation in trade series & 1945--90 & Trade, finance and insurance & $+$59 & $+$2 \\
Unknown Soviet military size & 1941--89 & Defence & $\pm$20 & $\pm$7 \\
No productivity adjustment in employment series & 1940--90 & Housing, education, health, science, administration & $-$147 & $-$10 \\
\bottomrule
\end{tabular}
\begin{tablenotes}[para,flushleft]
\footnotesize
\textit{Notes:} `Bias' is the assumed maximum error in the affected indicator, taken as 10 per cent in either direction unless an independent estimate exists; `Max.\ impact' is that error weighted by the affected sectors' share of GDP in the year in which it is largest. A minus sign means that growth in the new series may be understated, so that backcast levels are too high, and vice versa. Fixed value-added weights: Swedish single- versus double-deflated accounts \parencites[19]{swedengdp_15602010}[148--53]{lobellschonkrantz2008}; hidden inflation: difference between official and estimated real inflation in Soviet retail trade \parencite[607]{hidden_inflation_sov}; product taxes: in 1980--90 the ratio of product taxes less subsidies to value-added was 0.5 per cent in Bulgaria and Ukraine, the only socialist economies with data \parencite{wdi2024}; service productivity: British service-sector labour productivity rose by 147 per cent between 1948 and 1990 \parencite{statisticsUKlabour}. \textit{Sources:} \textcite[149--56]{klimantas2025thesis}; this paper.
\end{tablenotes}
\end{threeparttable}
\end{table}


Two deserve comment. The first is the index-number problem of a single late benchmark. A late base assigns to each sector the relative prices of the end of the period, and since the sectors that grew fastest were those whose relative prices fell, late weights understate growth; this is the `Gerschenkron effect', visible in Bergson's GNP for 1928--37, which grows by 175 per cent at 1928 prices and by 62 per cent at 1937 prices \parencite[517]{feinsteinthomas2002}. It is far smaller here. The effect is concentrated in the early phases of industrialisation \parencite[85]{jonassardy1970}, and for the post-war Soviet economy it is modest: the CIA's move from 1970 to 1982 prices lowered measured growth over 1951--84 by about half a point a year \parencites[6--7]{cia1990}[156]{harrison1993}. The bound comes from the data: if 1990 weights understated Latvian growth over 1940--90, the Soviet-period series would arrive in 1937--40 systematically above the 1935-weighted interwar series, and it does not. And the procedure is that of every comparable reconstruction, the CIA's 1982 weights for 1950--90 among them, in which Maddison found `nothing untoward' \parencite[312]{Maddison1998b}. The second is the fixed ratio of value-added to gross output. Thomas and Feinstein's simulations found that the procedure understated value-added growth by between 0.15 and 1 percentage point a year; the Swedish accounts, which allow single- and double-deflated series to be compared, show an understatement of the deflator of 17 per cent over the three-quarters of a century after 1875 \parencites[19]{swedengdp_15602010}[148--53]{lobellschonkrantz2008}, and the bound in Table~\ref{tab:sensitivity}, 14 per cent over the half-century, is in the lower part of that range, as befits indices that track output product by product and an economy that, like those of Central and Eastern Europe, could not renew the technological bases of its old branches of industry \parencite[195--7]{berend_centraleastern19441993}. The sign of the bias depends on the movement of input relative to output prices, and the Swedish figure is taken as a bound on its magnitude rather than as evidence of its direction in a planned economy. The error accumulates over the extrapolation and is largest in the earliest years of the Soviet-period series, precisely where the overlap shows the series within 3.5 per cent of an independent level.

Of the remaining biases, the largest in Latvia is defence. The Soviet garrison of 80,000 in 1990 was, in proportion to population, the heaviest in the Baltic republics, and the World Bank still counted some 300,000 military personnel and dependants, 11 per cent of the population, at the time of its 1993 report \parencite[157]{wbtransitionlatv}; the garrison's component of public services is interpolated between the 1940 and 1990 totals in step with public-administration employment, with a maximum bias of 7 per cent of GDP in either direction. The trade index uses turnover at comparable prices, the one value series in the reconstruction; Howard's estimates of hidden inflation in Soviet retail trade imply an upward bias of at most 2 per cent of GDP \parencites[607]{hidden_inflation_sov}[688--9]{kontorovich_2001}; in Latvia the only published price index, that of the kolkhoz markets in four cities, rose by 4 per cent a year in 1970--5 and 13 per cent a year in 1975--7, while the official indices recorded none \parencite[221]{years_of_dependence}. The employment-based indices for the public services make no allowance for productivity change, which understates service growth if productivity rose, with a maximum effect of 10 per cent of the level. Unregistered activity, put at 10--20 per cent of official GDP in 1989 in the IMF's alternative assumptions for Russia, was of a size that OECD economies also displayed \parencites[5]{baltictransitioncontraction}[90--1]{shadow_europe}. The population of the Abrene district at its transfer in 1944 is not counted in the sources available to us: King speaks only of `several thousands' of Latvian citizens, and Misiunas and Taagepera, like Hiden and Salmon, of about 2 per cent of the pre-war population \parencites[84]{king_1965_economic}[74]{years_of_dependence}[129--30]{hiden_salmon}; the share of about 2 per cent used in the population series is a rounded figure rather than a count, but under the Maddison Project's convention the border change has no bearing on income per head at any plausible value. Taken together, the quantifiable biases imply a margin of about 4 per cent around the interwar levels if they are independent, and of 8 per cent if they all point one way, and of 14 per cent for the largest single bias around the Soviet-period levels, or nearly a quarter in the earliest years if the two same-signed biases were both at their maxima; the overlap shows that the realised error is far smaller. The sub-period growth rates on which the argument of the following sections rests are robust to all of them. The overlap bounds only the net error over 1940--90, but the arithmetic is simple: even if the largest same-signed bias, a quarter of the level, were concentrated entirely in 1968--89, it would add about one point a year to a rate of \gr{1968}{1989} per cent, against \gr{1950}{1968} per cent in 1950--68. No plausible combination of the biases in Table~\ref{tab:sensitivity} would raise Latvian growth in 1968--89 to the level of 1950--68, reverse Latvia's position relative to the Soviet average, or alter its ranking among the three Baltic economies in 1938 or 1989.
